\chapter{Tổng quan}
\label{ch:tongquan}

Chương này dựng nền khái niệm cho phần còn lại của luận văn: kiến trúc mô hình được dùng,
kỹ thuật tinh chỉnh trong điều kiện tài nguyên hạn chế, bài toán định vị phần tử giao
diện, bộ dữ liệu nguồn, và sau cùng là vì sao các thước đo có
sẵn không dùng được cho bài toán này, kèm số đo cho từng cách.

\section{Từ học sâu tới mô hình thị giác--ngôn ngữ}

Học sâu là nhánh học máy dùng mạng nơ-ron nhiều tầng để tự học biểu diễn từ dữ liệu thô.
Bước ngoặt gần nhất với luận văn là kiến trúc Transformer với cơ chế tự chú
ý~\cite{vaswani2017}, cho phép mô hình ước lượng quan hệ giữa các đơn vị trong chuỗi mà
không phụ thuộc khoảng cách vị trí, đồng thời tận dụng được tính toán song song để mở
rộng quy mô tham số.

Mô hình thị giác--ngôn ngữ (\emph{vision--language model}, VLM) ghép một bộ mã hoá ảnh
vào một mô hình ngôn ngữ, thường qua một tầng chiếu (\emph{projector}) đưa đặc trưng ảnh
về cùng không gian với nhúng từ. Ảnh được cắt thành các mảnh, mã hoá thành một dãy
\emph{token thị giác}, rồi nối vào chuỗi văn bản; từ đó mô hình sinh văn bản có điều kiện
trên cả hai phương thức.

Mô hình gốc được chọn là \textbf{Qwen2.5-VL-3B-Instruct}~\cite{qwen25vl}, vì ba lý do.
Thứ nhất, đây là mô hình mở, tải được trọng số và tinh chỉnh được, khác với các mô hình
thương mại chỉ gọi qua giao diện lập trình. Thứ hai, ba tỉ tham số là cỡ chạy được tại
máy người dùng, phù hợp với động lực ứng dụng đã nêu ở Mục 1.2.2. Thứ ba, họ Qwen-VL xử
lý ảnh ở độ phân giải động, không ép ảnh về một kích thước cố định, điều có ý nghĩa
với ảnh màn hình điện thoại vốn có tỉ lệ rất dài, ví dụ $1080 \times 2400$.

\section{Tinh chỉnh hiệu quả tham số: LoRA và QLoRA}

Tinh chỉnh toàn bộ ba tỉ tham số đòi hỏi bộ nhớ vượt xa một card đồ hoạ đơn, vì ngoài
trọng số còn phải giữ gradient và trạng thái của bộ tối ưu. LoRA~\cite{lora} giải bằng
cách đóng băng trọng số gốc $W \in \mathbb{R}^{d \times k}$ và chỉ học một số hạng bổ
sung hạng thấp:
\begin{equation}
W' = W + \frac{\alpha}{r} BA, \qquad
B \in \mathbb{R}^{d \times r},\; A \in \mathbb{R}^{r \times k},\; r \ll \min(d,k),
\label{eq:lora}
\end{equation}
trong đó $r$ là hạng và $\alpha$ là hệ số tỉ lệ. Số tham số phải học giảm từ $dk$ xuống
$r(d+k)$.

QLoRA~\cite{qlora} đi thêm một bước: lượng tử hoá trọng số đóng băng xuống $4$ bit theo
định dạng NormalFloat (NF4) rồi mới gắn các bộ chuyển đổi LoRA lên trên. Trọng số gốc chỉ
được giải nén khi cần cho phép nhân, còn gradient chỉ chảy qua phần $A$, $B$ ở độ chính
xác cao.

Trong luận văn, cấu hình cụ thể (Mục~\ref{sec:cauhinh}) cho $14.966.784$ tham số học
được, tức $0{,}48\%$ số tham số của phần mô hình ngôn ngữ. Con số này được tính lại bằng tay từ bảy mô-đun
trên $36$ tầng và dùng làm phép kiểm cơ học trước mỗi lượt huấn luyện: nếu số tham số
không khớp thì cấu hình đã trượt ở đâu đó, chẳng hạn tháp thị giác không thực sự bị đóng
băng.

\section{Định vị phần tử giao diện}

Bài toán định vị phần tử giao diện (\emph{GUI grounding}) nhận một ảnh màn hình và một
biểu thức văn bản, trả về vị trí của phần tử được mô tả, thường ở dạng một điểm hoặc một
khung bao. Các mô hình tiêu biểu gồm SeeClick~\cite{seeclick},
OS-Atlas~\cite{osatlas} và UGround~\cite{uground}.

Quy ước chấm phổ biến, kế thừa từ AndroidInTheWild~\cite{aitw}, là ngưỡng dung sai
$\tau = 0{,}14$ tính theo từng trục: dự đoán được coi là trúng khi
$|p_x - g_x| \le \tau W$ và $|p_y - g_y| \le \tau H$, với $g$ là toạ độ đáp án và $W$,
$H$ là bề ngang và chiều cao màn hình. Cần chú ý vùng chấp nhận ở đây là một
\emph{hình chữ nhật} tính theo từng trục chứ không phải một đĩa tròn: trên màn
$1080 \times 2400$ nó tương ứng $\pm151$ điểm ảnh theo chiều ngang nhưng tới $\pm336$
điểm ảnh theo chiều dọc. Mục~\ref{sec:luattrung} cho thấy quy ước này quá lỏng khi
đối tượng được chấm là câu chữ, bởi nó vốn được đặt ra để chấm một bộ định vị.

Trong luận văn, một mô hình thuộc nhóm này giữ vai trò khác hẳn: nó giữ vai \textbf{dụng cụ
đo} chứ bản thân nó không được đem đo. Ba điều kiện được đặt ra cho nó ở
Mục~\ref{sec:botro}, và mọi điểm số trong luận văn đều là điểm số \emph{có điều kiện}
trên dụng cụ này.

\section{Bộ dữ liệu AndroidControl}

AndroidControl~\cite{androidcontrol} do nhóm Google DeepMind công bố tại NeurIPS 2024,
giấy phép CC0, gồm $15.283$ chuỗi thao tác do người thật thực hiện trên thiết bị Pixel
trong khoảng một năm. Mỗi bước lưu ảnh màn hình, loại thao tác, toạ độ chạm nếu có, và
trường quan trọng nhất đối với luận văn là \code{step\_instructions}, câu hướng dẫn do
chính người thao tác viết cho bước đó.

\begin{quote}
\textbf{Chỗ định vị của luận văn.} Trong bài gốc AndroidControl cũng như trong các công
trình dùng bộ này mà chúng tôi khảo sát (Aguvis~\cite{aguvis}, SeeClick~\cite{seeclick}), câu hướng dẫn từng
bước nằm ở \textbf{đầu vào}, dùng để gợi ý cho tác tử suy ra thao tác cần thực hiện. Ở
luận văn này, chính câu đó là \textbf{đích sinh} và là đầu ra duy nhất đem chấm. Hướng
sử dụng dữ liệu bị đảo ngược.
\end{quote}

Một tính chất của bộ dữ liệu cần nêu ngay vì nó chi phối trần của mọi kết quả: câu do
người viết rất cụt, trung vị chỉ khoảng sáu từ, dạng ``click on OK'' hay ``Select the
search result''. Câu càng cụt thì càng khó lần ra phần tử, và Mục~\ref{sec:tran} cho
thấy điều này kéo trần đo được xuống còn $75{,}7\%$.

\section{Vì sao các thước đo có sẵn không dùng được}
\label{sec:thuoc-gay}

Bài toán có sẵn câu chuẩn do người viết, nên phản xạ tự nhiên là so câu sinh ra với câu
chuẩn. Ba nhóm cách đã được thử, và mức độ thất bại của từng nhóm được đo lại; kết quả ở
Bảng~\ref{tab:thuocgay}.

\begin{table}[t]
\caption{Ba cách chấm dựa trên câu chuẩn, và mức độ thất bại đo được trên dữ liệu thật của
bài toán.}
\label{tab:thuocgay}
\centering\small
\renewcommand{\arraystretch}{1.25}
\begin{tabular}{@{}p{3.1cm}p{3.4cm}p{6.4cm}@{}}
\toprule
Cách chấm & Kết quả đo & Nguyên nhân thất bại \\
\midrule
So chuỗi &
bác nhầm $97{,}5\%$ số câu diễn đạt khác &
``search bar'' và ``magnifying glass'' không chung từ nào \\[0.4em]

So vector ngữ nghĩa &
AUC $0{,}336$ trên $51$ ca viết tay &
đo độ gần \emph{chủ đề}: cặp \code{inbox}--\code{outbox} được $0{,}76$, cao hơn cặp
\code{search}--\code{magnifying glass} chỉ $0{,}55$ \\[0.4em]

Bơm lỗi bản đầu &
AUC $1{,}000$ &
\textbf{hằng đẳng thức, không phải phép đo}: bộ bơm lỗi gọi chính hàm khớp của thước rồi
dán lại đầu ra, nên hai phân phối tách nhau theo định nghĩa \\
\bottomrule
\end{tabular}
\end{table}

Dòng thứ hai đáng phân tích kỹ vì nó phá vỡ một trực giác phổ biến. Nhúng ngữ nghĩa
được huấn luyện để đặt gần nhau những đoạn văn nói về cùng chủ đề. ``Hộp thư đến'' và
``hộp thư đi'' cùng nói về thư từ nên nằm gần nhau, dù chúng là \emph{hai nút khác nhau}
trên màn hình. Ngược lại, ``thanh tìm kiếm'' và ``kính lúp'' nói về hai chuyện khác nhau
(một là ô nhập, một là hình vẽ) dù chúng là \emph{cùng một nút}. Tín hiệu mà thước
cần là tính đồng nhất của phần tử, còn tín hiệu mà nhúng cung cấp là độ gần chủ đề. Hai
tín hiệu này không xếp cùng chiều, nên không tồn tại ngưỡng nào tách được chúng.

Dòng thứ ba ghi lại một sai lầm về phương pháp. Một bộ
kiểm chứng thước cho AUC bằng đúng $1{,}000$, giá trị không thể xuất hiện từ một phép đo trên dữ liệu thật. Truy lại thì
nhánh sinh câu ``diễn đạt khác'' được dựng bằng cách gọi chính hàm trích đích của thước
rồi ghép lại đầu ra, còn nhánh ``sai đích'' bị ép giao rỗng ở mức từ. Hai phân phối vì
vậy tách nhau theo định nghĩa, phương sai bằng không, và cổng chặn ``AUC dưới $0{,}80$
thì dừng'' không có đường nào để kích hoạt. Đo lại bằng các cặp diễn đạt thật cho AUC
$0{,}35$. Phép kiểm vì vậy được dựng lại với một ràng buộc bổ sung: bộ sinh dữ liệu kiểm không
được gọi hàm khớp của chính thước.

Kết luận của chương: câu hỏi ``câu sinh ra có giống câu chuẩn không'' không trả lời được
một cách đáng tin. Chương~\ref{ch:thuocdo} thay nó bằng một câu hỏi khác, ``câu sinh
ra có dẫn được tới đúng phần tử không'', và câu hỏi đó có barem khách quan là toạ độ
mà người dùng thật đã chạm.
